\section{来源、去重与可复核范围}\label{sec:provenance}
本文从包含420个 TPC 目录的独立仓库中选择核心原始来源，按数学依赖重写，不按论文编号串联。
正文中给出证明的命题仍以各自所列假设为准；经典 CRT、素数定理和 Mertens 输入不是本稿的新发现。
历史原稿的数量不等于独立贡献的数量。

配套文件 \texttt{SOURCE\_CATALOG.md/json} 枚举420个目录，保留原名称、主源文件散列和文本匹配的候选岛屿标签。
候选标签用于检索，不是对420篇全文逐篇进行人工语义审核的结论。
\texttt{CORE\_SOURCE\_MATRIX.md/json} 则记录本稿实际选用的核心文件、来源行区间、SHA-256 与声明边界。
源证明、历史接口、当前交接状态和独立审阅在其中分别列出，不以状态文件替代证明。

第~\ref{sec:structure}节的较早结构接口来自原始 TeX；
第~\ref{sec:operator}至第~\ref{sec:boundary}节的主要算子结果来自2026年9月9日未编号研究循环；
第~\ref{sec:dynamics}节的 CRT--Abel 证明来自移动靶原始证明，复制删除结构来自精确物理交织文件。
原有论文、研究原稿、交接文件和地图均保持不变。新稿单独存放，未新增 TPC 编号。

正文的可复核核心是共同对象的完整定义、逐步证明与来源散列。
编译成功、引文解析成功、字体嵌入和逐页视觉检查只验证文档呈现，不能验证定理。
独立审稿也只对其明确读取范围给出意见，不构成整个420目录的认证或新颖性证明。
